\scanpage{49}
\mspage{8}
% L1–7
\uline{NB} La donnée d'une base $(\alpha_i)_{i\in I}$ de $\Delta$ équivaut à la donnée d'un isom
\[
\Delta\simeq\underbrace{I\amalg I\amalg\unsure{\mathfrak P_2|\mathscr P_2}(I)}_{\substack{\text{avec la structure}\\\text{de graphe déjà}\\\text{explicitée}}}
\]
On a donc prouvé que tt autom de $S$ invariant \deleted{cette} la partie basique sous-jacente est induit (de façon unique), par \deleted{\unsure{trans}} transport de structure, d'une permutation de $I$.

% L8–16
Base duale $\alpha'$ d'une base $\alpha$ ; notations $\eta(\alpha),\unsure{\zeta|\xi}(\alpha)$

\uline{Prop} $\exists!\,\lambda_\Delta=\lambda\in\widehat E(\Delta)$ tel que $\boxed{\lambda\cdot s=1\qquad\forall s\in\Delta}$

On a \deleted{$\sum_{s\in S}s=9\lambda$} (pour tte base $\alpha=(\alpha_i)_{i\in I}$)
\[
\left\{\begin{aligned}
\lambda&=3\eta(\alpha)-\zeta(\alpha)&&1^\circ)\\
\eta(\alpha)+\underset{\substack{\shortparallel\\\eta'(\alpha)}}{\eta(\alpha')}&=2\lambda&&2^\circ)\\
\zeta(\alpha)+\underset{\substack{\shortparallel\\\zeta'(\alpha)}}{\zeta(\alpha')}&=4\lambda&&3^\circ)
\end{aligned}\right.
\]
d'où
\[
\sum_{i,j\in I}\zeta_{ij}(\alpha)=5\lambda\tag*{$4^\circ)$}
\]
\[
\sum_{s\in S}\deleted{\unsure{\zeta}}s=9\lambda\tag*{$5^\circ)$}
\]
% L17–22
L'unicité est claire, puisque les $s$ \inserted{$\in S$} engendrent $\widehat E(S)$ et que la forme int. est non dég. Pour l'existence, et les \inserted{autres} formules écrites (nb aussi bien la première, que \deleted{l'une} $2^\circ$ ou $3^\circ$ jointe à $4^\circ$ ou $5^\circ$, \deleted{\unsure{et}} \inserted{que} $4^\circ$ et $5^\circ$, caractérisent $\lambda$\unsure{]|)} l'existence est \deleted{évidente sur} \inserted{connue dans} le cas standard, donc évidente.

% L23–27
\uline{Corollaire} \inserted{L'élément $\lambda$ est invariant par autom.} \deleted{\unsure{(}} Le produit scalaire par $\lambda$ donne une suite exacte invariante par autom,
\diagram{\begin{tikzcd}[ampersand replacement=\&,column sep=small]
0\arrow[r] \& E\arrow[r] \& \widehat E\arrow[r] \& \mathbb Z\arrow[r] \&0
\end{tikzcd}}
où $E$ est l'\unsure{orthogonale|orthogonal} de $\lambda$

Donc la formation de cet orthogonal
% archive sheet number: 48
